% !PRINTED 247
\weberpage{247}
%<<CONT
\weq{4}{27D = 4(a^2+12c)^3 - (2a^3-72ac+27b^2)^2}
bestimmt ist, ist zugleich die Discriminante der biquadratischen Gleichung (1), und die Vorzeichen der Grössen $u$, $v$, $w$ sind durch die Bedingung
\weq{5}{uvw = -b}
beschränkt [\S.\,37, (3)].

Wenn die Discriminante negativ ist, so hat die Gleichung (3) ebenso wie (1) zwei conjugirt imaginäre Wurzeln.

Ist $D$ positiv und alle vier Wurzeln $\alpha$, $\beta$, $\gamma$, $\delta$ reell, so werden auch $u$, $v$, $w$ reell und also ihre Quadrate, d.\,h. die Wurzeln von (3), \gesperrt{positiv}.

Sind alle vier Wurzeln imaginär, etwa $\alpha$ mit $\beta$ und $\gamma$ mit $\delta$ conjugirt, so folgt aus (2), dass $v$ und $w$ rein imaginär, $u$ reell ist; also hat in diesem Falle die Gleichung (3) eine positive und zwei negative Wurzeln.

In beiden Fällen kann aber auch, wenn $b$ verschwindet, eine der Grössen $u$, $v$, $w$ gleich Null sein.

Mit Rücksicht auf die Ergebnisse des letzten Paragraphen über die cubische Gleichung kommen wir also zu folgendem Resultat:

\gesperrt{Die nothwendige und hinreichende Bedingung für die Existenz von vier verschiedenen reellen Wurzeln ist}
\weq{6}{D > 0, \quad a < 0, \quad a^2 - 4c > 0.}

In allen anderen Fällen, wo $D$ positiv ist, hat die Gleichung vier imaginäre Wurzeln.

Wir können auch für $D=0$, also im Falle der Gleichheit zweier Wurzeln, die Discussion vollständig durchführen.

Nehmen wir $\gamma = \delta$ an, so folgt aus (2), da wir die Annahme $\alpha+\beta+\gamma+\delta = 0$ gemacht haben,
\[
2v = 2w = \alpha - \beta
\]
\[
2u = \alpha+\beta-2\gamma = 2(\alpha+\beta)
\]
und daraus nach (3)
\weq{7}{\begin{aligned} -4a &= 2(\alpha+\beta)^2 + (\alpha-\beta)^2\\ 16(a^2-4c) &= (\alpha-\beta)^2[8(\alpha+\beta)^2+(\alpha-\beta)^2]. \end{aligned}}
Um zunächst den Fall zu erledigen, dass auch $\alpha=\beta$ ist, so erhalten wir aus (7) dafür die Bedingung $a^2-4c=0$, und die erste Gleichung (7) zeigt, dass $a$ negativ ist, wenn $\alpha$ und $\gamma$
%CONT>>
